
\chapter{PLANOS Y LÁMINAS DE DETALLE}

\begin{longtable}{L{1.6cm}L{5.5cm}L{4.5cm}L{2cm}}
	\caption{Relación de planos presentados}
	\label{tab:planos}\\
	
	\toprule
	\textbf{Código} & \textbf{Descripción} & \textbf{Archivo} & \textbf{Estado} \\
	\midrule
	\endfirsthead
	
	\toprule
	\textbf{Código} & \textbf{Descripción} & \textbf{Archivo} & \textbf{Estado} \\
	\midrule
	\endhead
	
	IE-02 &
	Plano de Instalaciones Eléctricas - Primer Piso &
	IE-02-primer-piso.png &
	Entregado \\
	
	IE-03 &
	Plano de Instalaciones Eléctricas - Segundo Piso &
	IE-03-segundo-piso.png &
	Entregado \\
	
	IE-04 &
	Plano de Instalaciones Eléctricas - Tercer Piso &
	IE-04-tercer-piso.png &
	Entregado \\
	
	IE-05 &
	Diagrama Unifilar y Cuadro de Cargas &
	IE-05-unifilar-cuadro-cargas.pdf &
	Entregado \\
	
	IE-06 &
	Detalle del Sistema de Puesta a Tierra &
	IE-06-puesta-tierra.pdf &
	Entregado \\
	
	\bottomrule
\end{longtable}

\section{Planos de Instalaciones Eléctricas}

A continuación se presentan los croquis arquitectónicos originales de la vivienda, sobre los cuales se ha desarrollado el diseño eléctrico. Estos planos constituyen la base arquitectónica del proyecto y reflejan la distribución real de ambientes por nivel.

\subsection{IE-02: Plano de Instalaciones Eléctricas - Primer Piso}

Se presenta el croquis correspondiente al primer nivel de la vivienda, donde se muestra la distribución de los ambientes: tienda, cocina, pasadizo y escalera.

\begin{figure}[H]
	\centering
	\includegraphics[
	width=\textwidth,
	height=0.76\textheight,
	keepaspectratio
	]{../fuentes/croquis/piso-1.png}
	\caption{Croquis arquitectónico - Primer Piso (base para diseño eléctrico)}
	\label{fig:piso1}
\end{figure}

\subsection{IE-03: Plano de Instalaciones Eléctricas - Segundo Piso}

Se presenta el croquis correspondiente al segundo nivel de la vivienda, donde se muestra la distribución de los ambientes: dormitorio principal, dormitorio 3, sala/hall, baño y escalera.

\begin{figure}[H]
	\centering
	\includegraphics[
	width=\textwidth,
	height=0.76\textheight,
	keepaspectratio
	]{../fuentes/croquis/piso-2.png}
	\caption{Croquis arquitectónico - Segundo Piso (base para diseño eléctrico)}
	\label{fig:piso2}
\end{figure}

\subsection{IE-04: Plano de Instalaciones Eléctricas - Tercer Piso}

Se presenta el croquis correspondiente al tercer nivel de la vivienda, donde se muestra la distribución de los ambientes: dormitorio 4, dormitorio 5, dormitorio 6, pasadizo, baño y escalera.

\begin{figure}[H]
	\centering
	\includegraphics[
	width=\textwidth,
	height=0.76\textheight,
	keepaspectratio
	]{../fuentes/croquis/piso-3.png}
	\caption{Croquis arquitectónico - Tercer Piso (base para diseño eléctrico)}
	\label{fig:piso3}
\end{figure}

\subsection{IE-05: Diagrama Unifilar y Cuadro de Cargas}

Se presenta el diagrama unifilar correspondiente a los circuitos derivados, tableros de distribución y sistema de puesta a tierra.

\begin{landscape}
	\begin{figure}[H]
		\centering
		\includegraphics[
		width=\textwidth,
		height=0.82\textheight,
		keepaspectratio
		]{../entregables/planos/pdf/IE-05-unifilar-cuadro-cargas.pdf}
		\caption{Diagrama Unifilar - Tablero General TG-01}
		\label{fig:unifilar}
	\end{figure}
\end{landscape}

\subsection{IE-06: Detalle del Sistema de Puesta a Tierra}

Se presenta la lámina de detalle correspondiente a la instalación del sistema de puesta a tierra (pozo a tierra).

\begin{landscape}
	\begin{figure}[H]
		\centering
		\includegraphics[
		width=\textwidth,
		height=0.82\textheight,
		keepaspectratio
		]{../entregables/planos/pdf/IE-06-puesta-tierra.pdf}
		\caption{Plano de Detalle del Sistema de Puesta a Tierra}
		\label{fig:puesta_tierra}
	\end{figure}
\end{landscape}

\section{Observaciones}

\begin{itemize}
	\item Los planos fueron elaborados considerando la distribución arquitectónica de la vivienda unifamiliar de tres niveles.
	
	\item La ubicación de luminarias, interruptores y tomacorrientes se realizó de acuerdo con los criterios de funcionalidad y seguridad establecidos en el diseño eléctrico residencial.
	
	\item Las canalizaciones eléctricas fueron trazadas procurando recorridos técnicamente adecuados y compatibles con la distribución de los ambientes.
	
	\item El diseño fue desarrollado considerando los criterios establecidos en el Código Nacional de Electricidad -- Utilización (CNE-U).
\end{itemize}
